\documentclass[a4paper,11pt]{article}

\usepackage{revy}
\usepackage[utf8]{inputenc}
\usepackage[T1]{fontenc}
\usepackage[danish]{babel}
\usepackage{amsmath,amssymb,amsfonts,amsthm,textcomp}


\revyname{MatematikRevyen}
\revyyear{2026}
% HUSK AT OPDATERE VERSIONSNUMMER
\version{0.1}
\eta{$1$ minutter}
\status{Færdig}

\title{Løvens A-lokale 3}
\author{Carl '21 \& Thais '22}

\begin{document}
\maketitle

\begin{roles}
\role{E}[Skuespiller] Studerende-entrepeneur (nervøs type)
\role{S}[Skuespiller] Speaker (meget få linjer, ku være i bandet?)
\role{L1}[Skuespiller] Forelæser-løve (Fabien)
\role{L2}[Skuespiller] Forelæser-løve (Ernst)
\role{L3}[Skuespiller] Forelæser-løve (Morten)
\end{roles}

\begin{props}
\item Bord med tre stole til løverne
\end{props}


\begin{sketch}
\says{F} Den næste iværksætter, vil pitche en ny måde at bevise sætninger på, hvor han... hvor.. det nok bedst han selv forklare det.

\says{E} Hej løver, jeg har en... har en helt ny... et helt nyt bevis

\scene{Intro-jingle}

\says{E} En dag da jeg prøvede at bevise, at der var virkelig mange primtal. Så tænkte jeg, hey! Hvad nu hvis der ikke var virkelig mange primtal? Og så fandt jeg ud af, at der ikke var virkelig ikke mange, ikke var primtal. Og det er kernen i mit koncept. Modstridsbeviser


\scene{Lys ned}


\end{sketch}
\end{document}
